\newpage
\pagenumbering{arabic}
\setcounter{page}{1}

\section{Introduction}

The rapid expansion of structured data over the past decade has fundamentally transformed how information is synthesized across scientific research, industrial applications, and public services~\cite{dimitrakis_survey_2020, zhang_review_2021}. From massive general-purpose repositories like Wikidata and DBpedia to highly specialized domain-specific knowledge bases, Knowledge Graphs (KGs) have become the central substrate for integrating heterogeneous data. However, the sheer heterogeneity of this data—characterized by diverse ontological structures, varying levels of granularity, and distinct naming conventions—presents a substantial barrier to effective retrieval for the average user~\cite{hoffner_survey_2017, diefenbach_core_2018}. Accessing these repositories typically requires mastering SPARQL, a complex formal query language, and possess deep familiarity with domain-specific vocabularies~\cite{pereira_systematic_2022}.
Knowledge Graphs (KGs) have become the indispensable backbone of information infrastructure across a multitude of critical sectors. In Healthcare and Life Sciences, cognitive systems such as IBM Watson have demonstrated how KGQA can assist medical professionals by linking clinical trial data, pharmaceutical research, and patient histories to suggest potential diagnoses. In the Energy and Industrial sectors, systems like PetroKG are used by major entities such as PetroChina to integrate multi-source seismic and drilling data for intelligent assistant decision-making, while Siemens Energy utilizes knowledge graphs to manage the complex logistics and maintenance of global gas turbine fleets. The Financial Services industry represents perhaps the most robust adoption of this technology; platforms like Quantexa and those deployed by Goldman Sachs leverage graph-based entity resolution for real-time fraud detection, anti-money laundering (AML) surveillance, and tracking complex ownership hierarchies. In Media and E-commerce, giants such as Netflix and Amazon utilize knowledge graphs to move beyond simple keyword matching, instead using "thematic similarities" and purchase ontologies to power personalized recommendation engines that account for over 80\% of their content consumption~\cite{antoniou_survey_2022}. Furthermore, in Tourism, systems such as the Hainan Tourism Knowledge Graph (TKG) or earlier research prototypes like LILOG and LD-SDS, enable users to pose complex queries about geographical entities, cultural landmarks, and seasonal travel logistics.
Despite the ubiquity of these applications, the "access gap" remains a primary bottleneck. The field of Text-to-SPARQL aims to remove this barrier by translating natural-language questions into executable SPARQL queries. Earlier surveys on semantic question answering and KGQA provide valuable foundations regarding classic pipelines and challenges, such as the lexical gap, entity ambiguity, and structural query complexity~\cite{hoffner_survey_2017, diefenbach_core_2018}. However, these foundational approaches largely predate the recent paradigm shift toward Large Language Models (LLMs) and do not systematize the emerging design space of prompting, retrieval augmentation, and agentic workflows~\cite{dimitrakis_survey_2020, antoniou_survey_2022}. Consequently, there is a clear need for a research framework that treats LLM-based Text-to-SPARQL as a first-class research area.
The emergence of Large Language Models (LLMs) has recently revolutionized the field of KGQA~\cite{dabramo_investigating_2025, meloni_exploring_2025} by enabling more natural and fluid interactions with structured data. However, despite their sophisticated language understanding, LLMs face a fundamental "topological challenge": they often struggle to generate queries that strictly align with the specific schema and graph structure of a target KG without external grounding. Traditional approaches to solving this—primarily through supervised fine-tuning—achieve respectable performance, with some benchmarks reaching 59\% precision. However, these "black-box" models are often domain-specific, requiring extensive retraining for every new knowledge graph and offering no transparency into the reasoning process that led to a specific query.
This thesis proposes a novel framework for constructing a domain-agnostic, LLM-driven agentic system for Text-to-SPARQL generation. By employing LLMs as intelligent orchestrators, this research aims to overcome the limitations of traditional fine-tuning while maintaining transparency through structured reasoning traces. Rather than relying on a monolithic generation step, our system coordinates multiple specialized agents—facilitating planning, linking, and query refinement—to ensure robust performance across heterogeneous knowledge graphs~\cite{perevalov_text--sparql_2025, zhao_agentict2srobust_2025}.
The foundation of the proposed solution lies in three interconnected components. First, an Agentic Schema Linking and Discovery system that utilizes Retrieval-Augmented Generation (RAG). Unlike traditional similarity-based approaches that rely on static mappings, this system is inherently dynamic. It leverages in-context learning with carefully selected demonstrations to reach high performance without the need for parameter updates~\cite{dabramo_investigating_2025}. By using RAG to discover graph structures at runtime, the model can generalize to previously unseen knowledge graphs in a zero-shot manner, eliminating the need for the costly and non-scalable task-specific pretraining that characterizes traditional baselines~\cite{brei_leveraging_2024}.
The foundation of the proposed solution lies in three interconnected components. First, an Agentic Schema Linking and Discovery system that utilizes Retrieval-Augmented Generation (RAG)~\cite{lewis2020rag}. Unlike traditional similarity-based approaches that rely on static mappings, this system is inherently dynamic. It uses RAG to discover graph structures at runtime, allowing the model to generalize to previously unseen knowledge graphs without requiring prior training or manual alignment with specific schema vocabularies.
Second, a Topological Grounding mechanism that enriches the generation process. By retrieving the immediate graph neighborhood—including incoming and outgoing edges—for candidate entities, the system provides the LLM with the necessary ontological context to resolve semantic ambiguity~\cite{ji_retrieval--reasoning_2026}. This ensures that candidate entities are grounded in their local graph topology, preserving the integrity of the knowledge graph's unique structure.
Third, an Iterative Execution-Guided Refinement loop. This component serves as a self-correction mechanism, automatically detecting and repairing queries that return syntax errors or empty result sets~\cite{zhao_agentict2srobust_2025}. By using real-time execution feedback to prompt the agent to "re-think" and refine its logic, the framework significantly boosts precision and system reliability.
The motivation for this research stems from a critical gap in the adaptability of current Semantic Web infrastructure. Despite the central role of Knowledge Graphs in modern AI, there remains a lack of comprehensive systems that can autonomously navigate, validate, and query heterogeneous data at scale without extensive, domain-specific training. Current practices often rely on "frozen" fine-tuned models that fail to adapt to evolving schemas or new data sources without significant manual intervention~\cite{mecharnia_performance_2025}. The inefficiencies associated with these methods become more pronounced as global databases grow in size and diversity. By proposing an LLM-driven, modular framework, this thesis seeks to develop a scalable solution that reduces dependency on manual curation and enhances the semantic alignment and accessibility of global data.
The primary objective of this study is to investigate the efficacy of LLM agents in navigating the complexity of Text-to-SPARQL translation. Specifically, this research aims to: (1) develop an LLM-driven agentic workflow that decomposes query generation into interpretable stages, (2) design a RAG-based schema linking strategy that enables cross-domain agnostic portability, (3) implement a self-correction loop that utilizes execution feedback to improve query accuracy, and (4) evaluate the framework on the QALD-10 benchmark~\cite{usbeck2019benchmarking} using multilingual Wikidata entities to demonstrate robustness across diverse languages and schemas.
The expected outcome is a framework that transforms natural language questions into semantically rich SPARQL queries across various domains. This thesis makes several key contributions. Methodologically, it introduces a comprehensive LLM-driven pipeline for zero-shot SPARQL generation. Empirically, it demonstrates that while the agentic workflow achieves a precision of 50\% on QALD-10—compared to 59\% for traditional fine-tuned models—it offers superior explainability and eliminates the need for task-specific pretraining. Practically, this work delivers a reproducible methodology published as open-source software, supporting broader goals of democratizing access to the Semantic Web and fostering data-driven discovery across complex knowledge systems.
The remainder of this thesis is organized as follows. Chapter 2 provides a comprehensive review of the evolution of Knowledge Graph Question Answering (KGQA), tracing the transition from classical semantic parsing to modern Large Language Model (LLM) architectures and agentic design patterns. Chapter 3 presents the complete methodology of the proposed framework, detailing the RAG-driven schema linking process, the multi-agent coordination logic, and the execution-guided self-correction loop. Chapter 4 describes the experimental setup, the benchmarking process on the QALD-10 dataset, and a comparative analysis of results. Chapter 5 discusses the broader implications of this research for the Semantic Web, specifically addressing the trade-offs between execution accuracy and domain-agnostic scalability, while also examining current limitations. Finally, Chapter 6 concludes with a summary of key findings and outlines directions for future work in cross-lingual and multi-graph query generation.
In summary, this thesis addresses a fundamental challenge in the Semantic Web: the accessibility of structured knowledge through intuitive natural language interfaces. By leveraging Large Language Models as intelligent orchestrators capable of runtime schema discovery and iterative refinement, this research contributes to the development of more accurate, transparent, and domain-agnostic Text-to-SPARQL systems. Ultimately, this work supports the democratization of knowledge graphs, paving the way for adaptable tools that enable data-driven discovery across diverse scientific and industrial domains without the need for specialized technical expertise.